% Cote 119: offprints, reports, courses, CV and letters (typed and printed part): archive cote 119, site page 209 (PDF page 210).
% typescript. Draft compared: none.
% AI-transcribed by Claude Opus 5.5 (claude-opus-5-5, adaptive thinking, effort high), Message Batches API, batch r119_full (msgbatch_01DwZNm5Y7YsKBWyAN4sFs8w), 2026-11-13; not yet checked by eye.
\scanpage{119/209}
\mspage{- 14 -}

\uline{Algèbre}

[11 bis] (en collaboration avec J. Dieudonné) Critères différentiels de régu-%
larité pour les localisés des algèbres analytiques, Journal of Algebra,
vol. 5, pp. 305-324 (1967).

\uline{Groupes algébriques}

[12 bis] Exposés 4 (Sur quelques propriétés fondamentales en théorie des inter-%
sections) et 5 (torsion homologique et sections rationnelles),in
Anneauxde Chow et applications, Sém. Chevalley à l'ENS, 1958, (36 p +
29 p.).

\uline{Groupes discrets}

[13 bis] Représentations linéaires et compactification profinie des groupes
discrets, Manuscripta Math. vol 2, pp. 375-396 (1970).

\uline{Groupes Formels}

[14 bis] Groupes de Barsotti-Tate et cristaux, Actes Congr. Int. math. 1970,
t. 1., pp. 431-436.

\uline{Géométrie Algébrique}

[15 bis] Techniques de descente et théorèmes d'existence en Géométrie Algébrique
(recueil des exposés Bourbaki n°\unsure{s|8} 182, 190, 195, 212, 221, 232, 236),
Secrétariat de l'IHP, rue Pierre Curie, Paris (1958-1962).

[16 bis] On the De Rham Cohomology of algebraic varieties, Publ. Math. IHES,
vol. 29, pp. 95-103 (1966).

[17 bis] Hodge's general conjecture is false for trivial reasons, Topology,
vol. 8, pp. 299-303 (1969).

[18 bis] Local Cohomology (notes by R. Hartshorne), Lecture Notes in Math.
n° 41 (1967), Springer.

[19 bis] (en coll. avec J.P. Murre) The tame fundamental group of a formal
neighbourhood .. . Lecture \unsure{Notes|Notés} in Math. n° 208 (1971), Springer.

[20 bis] (avec H. Seydi), Platitude d'une adhérence schématique et lemme de
Hironaka généralisé, Manuscripta Math.5, pp. 323-339 (1971).

\mspage{\hand{209}}
